% TOP_PROBLEM: {"id": 3211, "title": "Circular choosability of planar graphs", "category_id": 3, "category": {"id": 3, "name": "graph_theory", "display_name": "Graph Theory", "description": "Problems involving graphs, networks, and their properties.", "slug": "graph-theory", "order_index": 3, "created_at": "2026-07-31T15:26:25.671Z"}, "status": "open", "problem_number": "OPG-37619", "set_id": 12}
% FIRST_POSTED: 2026-10-01
% ENTRY_KIND: statement_only
% UnsolvedMath ID 3211; source catalogue code OPG-37619
% !TeX program = lualatex
% Definition and sources only; this file is not an LLM solution attempt.
\documentclass[11pt,a4paper]{article}
\usepackage[margin=25mm]{geometry}
\usepackage{fontspec}
\setmainfont{lmroman10-regular.otf}[BoldFont=lmroman10-bold.otf,
 ItalicFont=lmroman10-italic.otf,BoldItalicFont=lmroman10-bolditalic.otf]
\usepackage{amsmath,amssymb,mathtools}
\usepackage{unicode-math}
\setmathfont{latinmodern-math.otf}
\AtBeginDocument{\let\setminus\smallsetminus}
\usepackage{xurl,hyperref}
\hypersetup{colorlinks=true,urlcolor=blue}
\setcounter{secnumdepth}{0}
\setlength{\parindent}{0pt}
\setlength{\parskip}{5pt}
\setlength{\emergencystretch}{3em}
\begin{document}
\section*{Circular choosability of planar graphs}\label{3211}
{\small UnsolvedMath ID 3211 · OPG-37619 · Graph Theory · OpenGarden}
\subsection{Definitions and mathematical statement}
For positive integers $p,q$, a circular $(p,q)$-coloring of a finite simple graph $G$ is a map $c:V(G)\to\{0,1,\ldots,p-1\}$ for which
\[
 q\le |c(u)-c(v)|\le p-q\qquad(uv\in E(G)).
\]
This says that the cyclic distance between adjacent colors is at least $q$. A list assignment for this palette gives each vertex a subset $L(v)\subseteq\{0,\ldots,p-1\}$; a list coloring must additionally satisfy $c(v)\in L(v)$.

For a real number $t\ge1$, call $G$ circularly $t$-choosable if, for every pair of positive integers $p,q$ and every list assignment with $|L(v)|\ge tq$ at every vertex, a circular $(p,q)$-list-coloring exists. Since list sizes are integral this means at least $\lceil tq\rceil$ colors. If the palette is too small to admit the specified lists, that pair imposes no condition. Define the circular choice number by
\[
 \operatorname{cch}(G)=\inf\{t\ge1:G\text{ is circularly }t\text{-choosable}\}.
\]
The problem is to determine the exact value of $\sup_G\operatorname{cch}(G)$ over all finite planar graphs. A solution requires a universal upper bound and matching examples or a sequence approaching that bound. Planarity refers to existence of a crossing-free drawing, and no particular embedding enters the lists. The quantifiers over palettes and individual lists are essential: this is more restrictive than ordinary circular colorability.
\subsection{Short English statement}
Determine the largest possible circular list-coloring requirement among all finite planar graphs.
\subsection{Sources}
\begin{itemize}
\item[S1] ULAM.AI and the UnsolvedMath Contributors, supplied problems.json, record 3211 (OPG-37619), OpenGarden. Catalogue identity and source transcription; CC BY 4.0.
\url{https://huggingface.co/datasets/ulamai/UnsolvedMath}.
\item[S2] Open Problem Garden, Circular choosability of planar graphs. Original problem and discussion.
\url{http://www.openproblemgarden.org/op/circular_choosability_of_planar_graphs}.
\end{itemize}
\end{document}
